\section{Experimental Setup}\label{sec:setup}

The experiments test three expectations from the design. A tracker that
already fits its detector should keep its accuracy. A tracker that does not
fit, because it has no continuation step or because the detector or its
score scale changed, should gain. Both should hold on trackers and data that
were not seen during design.

\label{sec:hosts}Nine published trackers serve as hosts: ByteTrack
\cite{zhang2022bytetrack}, BoT-SORT \cite{aharon2022botsort}, OC-SORT
\cite{cao2023ocsort}, SparseTrack \cite{liu2025sparsetrack}, BoostTrack
\cite{stanojevic2024boosttrack}, PD-SORT \cite{wang2025pdsort},
Hybrid-SORT \cite{yang2024hybridsort}, C-TWiX \cite{miah2025twix}, and
TrackTrack \cite{shim2025tracktrack}. Each host uses the authors' code and
published configuration, and its contract is given in
Table~\ref{tab:contracts}. A tenth system, TOPICTrack
\cite{cao2025topictrack}, was selected under the post-freeze protocol, but
its reproduction failed (67.54 against the authors' 69.6 HOTA on MOT17
val-half), so it is excluded from the analysis.

In every comparison the baseline is the host alone: the published tracker,
run by us with its published configuration and the same detections. The
second run adds the frozen \scl{} and changes nothing else. Both runs use the
same environment and the same detections. Where the authors report their own
numbers, these are listed next to our reproduction of the host. All tracking
runs used a CPU only, without a graphics processing unit (GPU).
Where this changed the numerical precision or the software environment of an
external system, the difference is given with its reproduction
(Sect.~\ref{sec:external}).
